% ========================================================================
% فصل ۹: مدیریت هویت و طرح‌های تصدیق هویت
% منابع: mazzocca2025survey، feraudoDIVADIDbasedReputation2024،
% luPseudonymChangingSocial2012، al-aniPrivacySafetyImprovement2023،
% fengBPASBlockchainAssistedPrivacyPreserving2020، mariaBAIVEfficientBlockchainBased2022،
% linGSISSecurePrivacyPreserving2007، rahmawatiagustinaSystematicLiteratureReview2025
% ========================================================================
\section{مدیریت هویت و طرح‌های تصدیق هویت}

\subsection{مدیریت هویت چیست و تکامل آن}

در ارتباطات شبکه اقتضایی خودرویی، هر پیام امنیتی باید به موجودیتی مشخص و قابل اعتماد پیوند خورد؛ از این رو شناسایی موجودیت‌ها و اثبات هویت آن‌ها پیش‌شرط هر سازوکار امنیتی است. هویت دیجیتال نیز تنها به انسان محدود نمی‌شود و دستگاه‌های اینترنت اشیاء در شبکه‌های نسل پنجم و ششم نیز به هویت یکتا نیاز دارند \cite{mazzocca2025survey}.

رویکردهای مدیریت هویت در طول چهار دوره تکامل یافته‌اند \cite{mazzocca2025survey}. در دورهٔ نخست، هویت متمرکز، هویت‌ها توسط یک مرجع مرکزی مدیریت می‌شدند؛ نخستین نمونه‌ها به سال ۱۹۸۸ بازمی‌گردد، وقتی که ادارهٔ اعطای اعداد اینترنت \LTRfootnote{\lr{Internet Assigned Numbers Authority (IANA)}} اعتبار آدرس‌های \lr{IP} را تعیین می‌کرد و بعدها ادارهٔ اعداد و نام‌های اینترنت \LTRfootnote{\lr{Internet Corporation for Assigned Names and Numbers (ICANN)}} برای نام‌های دامنه ایفا گردید \cite{mazzocca2025survey}. این مدل با مشکلاتی همراه بود: بهره‌گیری از نام کاربری و گذرواژه و آسیب‌پذیری آن‌ها در برابر حملات واژه‌نامه‌ای و جعل، تکه‌تکه شدن هویت به‌گونه‌ای که کاربر به اندازهٔ تعداد سرویس‌ها هویت‌های متعددی داشت، وجود نقطهٔ شکست منفرد و احتمال افشای داده، و هزینهٔ سنگین نگهداری داده برای ارائه‌دهندگان \cite{mazzocca2025survey}. در دورهٔ دوم، هویت فدرال، ورود یکپارچه به سرویس‌های مختلف \LTRfootnote{\lr{Single Sign-On (SSO)}} امکان‌پذیر شد؛ مایکروسافت پاسپورت پیشگام این دوره بود، در سال ۲۰۰۱ شرکت سان مایکروسیستمز اتحاد لیبرتی \LTRfootnote{\lr{Liberty Alliance}} را تشکیل داد و امروزه گوگل و متا این نقش را ایفا می‌کنند، اما ارائه‌دهندگان هویت همچنان متمرکزند و کنترل کاربر محدود است \cite{mazzocca2025survey}. در دورهٔ سوم، هویت کاربرمحور، کاربران بدون فدراسیون متمرکز هویت خود را مدیریت می‌کنند؛ در سال ۲۰۰۵ ابزارهای مشترکات هویت \LTRfootnote{\lr{Identity Commons}} نقش پررنگی در ترویج کارگاه هویت اینترنت \LTRfootnote{\lr{Internet Identity Workshop (IIW)}} داشتند و پروژه‌هایی مانند \lr{OpenID}، نسخهٔ ۲.۰ آن، \lr{OIDC}، \lr{OAuth} و \lr{FIDO} مطرح شدند، اما همچنان به اعتماد به ارائه‌دهندگان سرویس وابسته‌اند \cite{mazzocca2025survey}. در دورهٔ چهارم، هویت خودمختار که در سال ۲۰۱۲ معرفی شد، فرد را در مرکز کنترل داده‌های خود قرار می‌دهد \cite{mazzocca2025survey}.

\subsection{هویت خودمختار \lr{(SSI)}}

هویت خودمختار \LTRfootnote{\lr{Self-Sovereign Identity (SSI)}} به‌عنوان دورهٔ پایانی این تکامل، فرد را در مرکز قرار می‌دهد و به او اجازه می‌دهد تصمیم بگیرد چه زمانی، آیا و چگونه داده‌های خود را افشا یا تغییر دهد \cite{mazzocca2025survey}. این رویکرد بر زیرساخت‌های غیرمتمرکز و فناوری‌هایی چون دفترهای کل توزیع‌شده، شناسه‌های غیرمتمرکز و اعتبارنامه‌های تأییدپذیر استوار است؛ با بهره‌گیری از دفتر کل توزیع‌شده‌ای مانند زنجیره بلوکی نیازی به مرجع‌های مرکزی نیست و داده‌ها رمزنگاری‌شده و قابل تأیید بوده و شفافیت و تغییرناپذیری را فراهم می‌کنند \cite{mazzocca2025survey}. در این سازوکار، هر فرد با شناسه‌ای غیرمتمرکز که به جفت‌کلیدی تحت کنترل خود متصل است شناسایی می‌شود، کلید عمومی معمولاً بر دفتر کل توزیع‌شده ثبت می‌شود، و ادعاها با اعتبارنامه‌های تأییدپذیر اثبات می‌شوند بدون آنکه نیازی به تعامل با مرجع صادرکنندهٔ متمرکز باشد؛ اعتبارنامه‌ها به‌دلایل حریم خصوصی معمولاً خارج از زنجیره منتقل می‌شوند \cite{mazzocca2025survey}. دستاورد بنیادین این رویکرد، ارائهٔ یک اعتبارنامهٔ معتبر به طرف سوم بدون واسطه است؛ نمونه‌ای از آن، استفاده از مدرک دانشگاهی امضا‌شده در فرآیند جست‌وجوی کار است \cite{mazzocca2025survey}.

انگیزه‌های مطرح شدن این رویکرد ریشه در ضعف‌های دوره‌های پیشین دارد: مخازن دادهٔ متمرکز می‌توانند به افشاهای جدی داده منجر شوند و سامانه‌ها در برابر مسئلهٔ نقطهٔ شکست منفرد آسیب‌پذیرند \cite{mazzocca2025survey}. در دورهٔ فدرال، کاربران کنترل محدودی بر داده‌های خود داشتند و مقرراتی مانند مقرریت حمایت از داده‌های عمومی اتحادیهٔ اروپا بر این اصل تأکید می‌کنند که فرد باید کنترل کامل اطلاعات خود را داشته باشد بدون آنکه ناچار باشد داده‌های خود را به مرجع متمرکزی یا طرف سومی واگذار کند \cite{mazzocca2025survey}. این جریان مقرراتی در مه ۲۰۲۴ با مقرریت ۲۰۲۴/۱۱۸۳ اتحادیهٔ اروپا و بنیان‌گذاری چارچوب هویت دیجیتال اروپایی تثبیت شد \cite{mazzocca2025survey}.

\begin{figure}[H]
\centering
\includegraphics[width=0.85\textwidth]{SSI vs centeralized.png}
\caption{مقایسه معماری مدیریت هویت متمرکز و هویت خودمختار}
\label{fig:ssi_vs_centralized}
\end{figure}

شکل \ref{fig:ssi_vs_centralized} معماری مدیریت هویت متمرکز را در برابر هویت خودمختار مقایسه می‌کند و نقش محوری فرد در مدل دوم را نشان می‌دهد.

\subsection{شناسه‌های غیرمتمرکز \lr{(DID)}}

پیاده‌سازی هویت خودمختار بر شناسه‌های غیرمتمرکز \LTRfootnote{\lr{Decentralized Identifiers (DID)}} استوار است که پس از تلاش‌های مشارکتی گسترده در سال‌های ۲۰۱۷ تا ۲۰۱۹، مشخصات آن به‌عنوان توصیهٔ رسمی \lr{W3C} منتشر شد \cite{mazzocca2025survey}. شناسهٔ غیرمتمرکز یک طرح شناسه‌گذاری جهانی یکتا و رمزنگاری‌شده است و موضوع شناسهٔ غیرمتمرکز — که می‌تواند انسان یا موجودیت غیرانسانی باشد — را به‌طور یکتا شناسایی می‌کند \cite{mazzocca2025survey}. از نظر نحوی، هر شناسه از سه بخش تشکیل شده است: شناسهٔ منبع یکنواخت، شناسهٔ مربوط به روش خاص شناسه، و شناسهٔ مخصوص آن روش؛ نمونه‌ای از آن \lr{did:example:123456789abcdefghi} است \cite{mazzocca2025survey}. روش شناسهٔ غیرمتمرکز فرآیندهای ایجاد، حل، به‌روزرسانی و غیرفعال‌سازی شناسه‌ها و اسناد شناسه را مشخص می‌کند و شناسهٔ یونیفرم منبع شناسه با افزودن مسیر، پرس‌وجو و تکه‌گوره گسترش می‌یابد تا منبعی را مکان‌یابی کند \cite{mazzocca2025survey}.

سند شناسهٔ غیرمتمرکز، سندی ماشین‌خوان به قالب \lr{JSON-LD} است که اطلاعاتی دربارهٔ موضوع شناسه دارد، از جمله کلیدهای عمومی رمزنگاری، نقاط اتصال سرویس، پارامترهای احراز هویت، مهرهای زمانی و متادیتای تکمیلی \cite{mazzocca2025survey}. این شناسه‌ها پایا و دائمی هستند و شناسایی قابل‌اعتمیدی را حتی با تغییر ارائه‌دهندهٔ سرویس فراهم می‌کنند \cite{mazzocca2025survey}. مالکیت شناسه با کلید خصوصی متناظر با کلید عمومی ثبت‌شده در سند اثبات می‌شود و تأییدکننده سند را از ثبت دادهٔ تأییدپذیر می‌خواند \cite{mazzocca2025survey}. کنترل‌گر شناسه، موجودیتی مجاز به تغییر سند شناسه است و ممکن است چند کنترل‌گر وجود داشته باشد یا کنترل‌گر با موضوع یکی باشد \cite{mazzocca2025survey}.

شناسه‌های غیرمتمرکز از جنبهٔ دامنهٔ کاربرد به سه دسته تقسیم می‌شوند که این طبقه‌بندی در مرور مذکور به مشخصات روش شناسهٔ هم‌سطح (Peer DID) منتسب شده و جزء مشخصات هستهٔ \lr{DID} نیست \cite{mazzocca2025survey}: شناسه‌های بی‌طرف \LTRfootnote{\lr{Anywise}} که با تعداد نامشخصی از طرف‌ها، معمولاً غریبه‌ها، قابل استفاده‌اند؛ شناسه‌های دوتایی \LTRfootnote{\lr{Pairwise}} که تنها توسط شخص و یک طرف دیگر شناخته می‌شوند و هر رابطه شناسهٔ یکتای خود را دارد و خطر همبستگی را به‌حداقل می‌رساند؛ و شناسه‌های چندطرفه \LTRfootnote{\lr{N-wise}} که به‌طور دقیق توسط \lr{N} طرف شامل شخص موضوع شناخته می‌شوند و شناسه‌های دوتایی را در حالت خاصی که \lr{N} برابر دو است، در بر می‌گیرند \cite{mazzocca2025survey}.

حل شناسه‌های غیرمتمرکز از طریق یک حل‌کنندهٔ جهانی انجام می‌شود که از سامانه‌های شناسهٔ متعدد پشتیبانی می‌کند و هر سامانه باید یک تطبیق‌دهندهٔ شناسه بین روش‌های خاص و حل‌کنندهٔ جهانی پیاده کند \cite{mazzocca2025survey}. در پروتکل احراز هویت شناسهٔ غیرمتمرکز، یک چرخهٔ چالش و پاسخ برای اثبات کنترل شناسه به‌کار می‌رود و این پروتکل جایگزین بالقوه‌ای برای نام کاربری و گذرواژه است؛ چارچوب \lr{DIDComm} نیز به‌عنوان بستر تعامل‌پذیری مطرح شده است \cite{mazzocca2025survey}.

\subsection{اعتبارنامه‌های تأییدپذیر \lr{(VC)}}

برای اثبات ادعا دربارهٔ موضوع شناسه، اعتبارنامه‌های تأییدپذیر \LTRfootnote{\lr{Verifiable Credentials (VC)}} به‌کار می‌رود که مشخصاتی استانداردشده توسط \lr{W3C} برای ساختن ساختار داده‌ای درهم‌تنیده قابلیت نمایش ادعاها — مانند خواص یا ویژگی‌ها — به‌گونه‌ای رمزنگاری‌شده، قابل تأیید و در برابر دستکاری مقاوم است \cite{mazzocca2025survey}. اعتبارنامه‌ها در کیف پول دیجیتال نگهداری و قابل‌حمل هستند و نمونه‌ای از آن‌ها اعتبارنامهٔ فارغ‌التحصیلی دانشگاه نمونه با اثبات از نوع \lr{RsaSignature2018} است \cite{mazzocca2025survey}. سه نقش اصلی در این سامانه هستند: نگهدارنده \LTRfootnote{\lr{Holder}} موجودیتی است که بر یک یا چند اعتبارنامه کنترل ممارسه می‌کند، صادرکننده \LTRfootnote{\lr{Issuer}} نهادهای مورد اعتمادی مانند نهادهای دولتی یا بانک‌هاست و تأییدکننده \LTRfootnote{\lr{Verifier}} موجودیتی است که برای ارائهٔ سرویس به اعتبارنامهٔ معتبر نیاز دارد \cite{mazzocca2025survey}. در گردش کار استاندارد، صادرکننده شناسه‌ها و طرح‌ها را در ثبت دادهٔ تأییدپذیر ثبت می‌کند، اعتبارنامه را به نگهدارنده صادر می‌کند، نگهدارنده ارائهٔ تأییدپذیر را برای تأییدکننده می‌فرستد و تأییدکننده شناسه‌ها و طرح‌ها را در ثبت بررسی می‌کند \cite{mazzocca2025survey}. ثبت دادهٔ تأییدپذیر \LTRfootnote{\lr{Verifiable Data Registry (VDR)}} میانجی مورد اعتمادی است که به‌عنوان مخزن اطلاعات اساسی از جمله اسناد شناسهٔ غیرمتمرکز عمل می‌کند و چرخهٔ حیات ایجاد، ثبت و ابطال شناسه‌ها و اعتبارنامه‌ها را مدیریت می‌کند؛ در حالی که اعتبارنامه‌ها معمولاً به‌جای نگهداری در ثبت، خارج از زنجیره به اشتراک گذاشته می‌شوند \cite{mazzocca2025survey}. ساختار اعتبارنامه شامل شناسهٔ منبع، شناسهٔ صادرکننده و شناسهٔ یکتای اعتبارنامه است که می‌تواند شناسهٔ غیرمتمرکز باشد و شرایط انقضای ادعا و امضاهای رمزنگاری را دربر می‌گیرد \cite{mazzocca2025survey}. ارائهٔ تأییدپذیر \LTRfootnote{\lr{Verifiable Presentation (VP)}} نیز روش‌های امضا کردن و ارائهٔ اعتبارنامه‌ها توسط نگهدارنده را مشخص می‌کند و قالب‌های آن می‌تواند \lr{JSON-LD}، \lr{JSON} یا نشانهٔ وب \lr{JSON} باشد؛ حضور چالش و دامنه در اثبات ارائه، به مقابله با حملهٔ پخش مجدد کمک می‌کند \cite{mazzocca2025survey}.

اعتبارنامه‌ها امکان افشای انتخابی بخشی از اطلاعات را دارند و دسته‌های آن شامل ادعاهای تکی، مقادیر هش‌شده، اثباتات دانش صفر و امضاهای افشای انتخابی است؛ راه‌حل روزافزون \lr{SD-JWT} با جایگزینی ادعاهای متن ساده با پیچش‌های مقادیر نمک‌شده، حریم خصوصی را تقویت می‌کند \cite{mazzocca2025survey}. نمونه‌هایی از آن، خرید شراب با افشای صرف تاریخ تولد یا اثبات بیمه بدون افشای سابقهٔ پزشکی است \cite{mazzocca2025survey}. از محدودیت‌های \lr{SD-JWT} این است که اندازهٔ اعتبارنامه به‌صورت خطی با تعداد ادعاها رشد می‌کند و تعداد دقیق ادعاها را افشا می‌کند و از این رو در معرض حملهٔ استنتاج است \cite{mazzocca2025survey}. اثبات دانش صفر \LTRfootnote{\lr{Zero-knowledge Proof (ZKP)}} به‌عنوان یکی از این دسته‌ها، توسط چارچوب هویت \lr{IOTA} با عنوان افشای انتخابی دانش صفر پشتیبانی می‌شود و در جدول مقایسهٔ پلتفرم‌های هویت، \lr{BBS+} و \lr{SD-JWT} و افشای انتخابی دانش صفر برای ابزار هویت هایپرلجر آریز نیز آمده است \cite{mazzocca2025survey}.

ابطال اعتبارنامه‌ها چالشی پایدار است؛ چون اعتبار اعتبارنامه در طول زمان تغییر می‌کند، نخستین راه‌حل تعیین دورهٔ اعتبار در خود اعتبارنامه است، اما برای ابطال در اثر سلب امتیاز، سامانه‌های بررسی وضعیت گواهی آنلاین و فهرست‌های ابطال گواهی از زیرساخت کلید عمومی و «فهرست ابطال ۲۰۲۰» \lr{W3C} به‌کار می‌رود که رشته‌ای بیتی است و با قرار گرفتن هر بیت روی ۱، اعتبارنامهٔ متناظر ابطال می‌شود و با فشرده‌سازی \lr{ZLIB} قابل فشرده‌سازی است \cite{mazzocca2025survey}. چالش اینجاست که در زمان انتشار مرور، تنها یک مشخصهٔ پیشنهادی وجود داشته که هنوز به رتبهٔ استاندارد \lr{W3C} نرسیده و ابطال کارا در اینترنت اشیاء و مقیاس‌پذیری آن در محیط‌های پرترافیک و دستگاه‌های محدود، هنوز کم‌بررسی است \cite{mazzocca2025survey}.

\subsection{مشکلات امنیتی رفع‌شده توسط \lr{SSI} و \lr{DID}}

با این تعریف‌ها، مشخص است که هویت خودمختار و شناسه‌های غیرمتمرکز به چالش‌های امنیتی شبکه‌های خودرویی پاسخ می‌دهند. در ارتباطات خودرو با خودرو، چالش کلیدی فقدان اعتماد میان خودروهاست، زیرا خودروها معمولاً رابطهٔ پیشینی ندارند و به مرجع‌های شبکهٔ متمرکز متکی‌اند؛ در این زمینه، شناسه‌های غیرمتمرکز و اعتبارنامه‌های تأییدپذیر یکپارچگی، احراز هویت، محرمانگی و حریم خصوصی را بدون اتکا به مرجع متمرکز مورد اعتماد تضمین می‌کنند \cite{mazzocca2025survey}.

نخستین دستاورد، حذف نقطهٔ شکست منفرد و کاهش وابستگی به مرجع متمرکز است \cite{mazzocca2025survey}؛ هرچند باید توجه داشت که در طرح‌هایی مانند \lr{DIVA}، مرجع مورد اعتماد همچنان شناسه‌های غیرمتمرکز را ثبت و اعتبارنامه صادر می‌کند و وابستگی کاهش یافته نه حذف‌شده است \cite{feraudoDIVADIDbasedReputation2024}. دوم، این فناوری‌ها حریم خصوصی را تقویت می‌کنند: در سامانهٔ \lr{DIVA} شناسه‌های غیرمتمرکز به‌عنوان شبه‌نام عمل می‌کنند و قابل نسبت دادن به هویت واقعی کاربر نیستند \cite{feraudoDIVADIDbasedReputation2024} و شناسه‌های دوتایی خطر همبستگی را کاهش می‌دهند، هرچند به‌طور کامل از پیوند جلوگیری نمی‌کنند؛ از این رو استفاده از شناسه‌های پایدار در اعتبارنامه‌ها توصیه نمی‌شود و چون تأییدکننده برای جلوگیری از حملهٔ سیبل باید ارائه‌های تکراری را بشناسد، رویکردهای «پیوندناپذیری پیوندشده» پیشنهاد شده است \cite{mazzocca2025survey}. سوم، افشای انتخابی اطلاعات بیش از حد را ممکن می‌سازد و رمزنگاری ادعاها برای تأییدکننده به‌عنوان دفاع در برابر افشای اطلاعات در نظر گرفته شده است \cite{mazzocca2025survey}. چهارم، در برابر تهدیدهای رمزنگاری، چرخش کلید، احراز هویت چندعاملی، ماژول‌های امنیتی سخت‌افزاری و بازیابی از طرف سوم برای خطر از دست دادن کلید، درج موضوع شناسه در اعتبارنامه و اثبات دانستن کلید خصوصی برای سرقت و تبانی اعتبارنامه، امضای صادرکننده و ثبت‌های صادرکنندگان مورد اعتماد برای جعل اعتبارنامه، و رمزنگاری سرتاسری برای حملهٔ فردی در میان و نشانه یک‌بارمصرف یا مهر زمانی و چالش برای پخش مجدد پیشنهاد می‌شوند \cite{mazzocca2025survey}.

پنجم، ردیابی شرطی در این معماری مسئله‌ای ظریف است. در \lr{DIVA}، شناسه‌های غیرمتمرکز به هویت واقعی کاربر نمی‌توانند نسبت داده شوند و ردیابی تنها از طریق فیلتر کردن شهرت امکان‌پذیر است که این یک شکاف پژوهشی محسوب می‌شود \cite{feraudoDIVADIDbasedReputation2024}. در مقابل، رویکردهایی برای حفظ شهرت هنگام تغییر آدرس بلاکچین پیشنهاد شده‌اند که در آن‌ها آدرس جدید از راز مشتق‌شده از کلید خصوصی شناسهٔ غیرمتمرکز ساخته می‌شود و این آدرس‌ها زنجیره‌ای قاطع می‌سازند که برای حسابرسی یا تحقیق قابل بازسازی است و ادعای آن‌ها با اثبات دانش صفر به قرارداد هوشمند و درخت مرکل تقویت می‌شود \cite{mazzocca2025survey}. ششم، مدیریت شهرت غیرمتمرزد نیز امکان‌پذیر است؛ در \lr{DIVA} امتیازات شهرت به‌صورت تراکنش‌های بدون ارزش در یک دفتر کل با دسترسی محدود ذخیره می‌شوند \cite{feraudoDIVADIDbasedReputation2024}.

در حوزهٔ حمل‌ونقل هوشمند نیز کاربردهای عملی این فناوری‌ها گزارش شده است. استاندارد شناسهٔ خودرو \lr{VID} مبتنی بر بلاکچین که در سال ۲۰۱۹ توسط ابتکار امنیت و تعامل متحرک مبتنی بر بلاکچین \LTRfootnote{\lr{Mobility Open Blockchain Initiative (MOBI)}} تعریف شد، شمارهٔ شناسایی خودرو را با بلاکچین سازگار می‌کند و معماری مرجعی با کیف پول خودرو شامل اعتبارنامه‌های مالکیت، تولد، ثبت، گارانتی، بیمه، هویت، میلیاژ و شناسه و صادرکنندگانی چون نهادهای دولتی، خودروساز و واحد کنارجاده‌ای ارائه می‌دهد \cite{mazzocca2025survey}. در طرح \lr{D-V2X}، واسطهٔ مورد اعتماد حذف شده و زیرساخت کلید عمومی غیرمتمرکز جایگزین زیرساخت سنتی می‌شود؛ خودرو شناسهٔ غیرمتمرکز تصادفی ثبت می‌کند، خودروساز اعتبارنامه‌ای حاوی شمارهٔ شناسایی خودرو صادر می‌کند و خودرو هم موضوع و هم نگهدارنده است؛ از منظر حریم خصوصی، خودرو چندین شناسهٔ غیرمتمرکز ثبت می‌کند که همه از طریق اعتبارنامه‌های مبتنی بر اثبات دانش صفر صادرشده توسط خودروساز به شناسهٔ اصلی پیوند دارند، یا از محاسبهٔ چندطرفه بهره می‌گیرد و همچنین می‌تواند با ساختن دو ارائهٔ تأییدپذیر از شبه‌نام استفاده کند \cite{mazzocca2025survey}. در برابر جعل \lr{GPS}، اعتبارنامهٔ مکان از زیرساخت نزدیک — مانند چراغ راهنمایی — استفاده می‌شود و رویکردهای اصالت داده با واحد کنارجاده‌ای و پیام‌های ایمنی پایه نیز گزارش شده است \cite{mazzocca2025survey}. سامانهٔ مدیریت کلید غیرمتمرکز در ارتباطات خودرو با همه‌چیز نیز اعتبارنامه‌ای پیونددهندهٔ شناسهٔ غیرمتمرکز به شمارهٔ شناسایی خودرو دارد \cite{mazzocca2025survey}. در نهایت، طرح ثبت و احراز هویت امن با بلاکچین دولایه و شناسه‌های غیرمتمرکز، بلاکچینی دولایه را معرفی می‌کند که در لایهٔ بالا واحدهای کنارجاده‌ای مجاز و در لایهٔ پایین واحدها و خودروهای پوشش قرار دارند؛ خودرو شناسهٔ غیرمتمرکز خود را می‌سازد، واحدهای کنارجاده‌ای اعتبارنامهٔ یکتا صادر می‌کنند، گیرنده شناسهٔ فرستنده را با فهرست مجاز بررسی و آستانهٔ شهرت را کنترل می‌کند و بازخورد را برای به‌روزرسانیٔ شهرت به واحد کنارجاده‌ای می‌فرستد \cite{mazzocca2025survey}.

\subsection{طرح‌های موجود برای تصدیق هویت}

در عمل، استفاده از شبه‌نام — که موضوع فصل ۵ بود — تنها بخشی از مسئله است؛ چالش اصلی تغییر شبه‌نام در زمان و مکان مناسب است، زیرا اگر خودرو شبه‌نام خود را در موقعیت نامناسبی تغییر دهد، تغییر شبه‌نام بی‌اثر خواهد بود \cite{luPseudonymChangingSocial2012, al-aniPrivacySafetyImprovement2023, rahmawatiagustinaSystematicLiteratureReview2025}. در این بخش، چهار طرح نمایان از طرح‌های تصدیق هویت و حفظ حریم خصوصی بررسی می‌شود: تغییر شبه‌نام در نقاط اجتماعی، طرح حریم خصوصی مرتبط با ایمنی، سامانهٔ شهرت مبتنی بر شناسهٔ غیرمتمرکز و پروتکل‌های مبتنی بر زنجیره بلوکی.

\subsubsection{تغییر شبه‌نام در نقاط اجتماعی \lr{(PCS)}}

رویکرد تغییر شبه‌نام در نقاط اجتماعی \LTRfootnote{\lr{Pseudonym Changing at Social Spots (PCS)}} پیشنهاد می‌کند که خودروها شبه‌نام خود را در نقاط اجتماعی تغییر دهند؛ نقاط اجتماعی مکان‌هایی هستند که چندین خودرو به‌طور موقت در آن‌ها جمع می‌شوند و به دو دستهٔ کوچک و بزرگ تقسیم می‌شوند: نقاط کوچک مانند تقاطع‌های چراغ قرمز و نقاط بزرگ مانند پارکینگ‌های رایگان نزدیک مرکز خرید \cite{luPseudonymChangingSocial2012}. الگوریتم این طرح ساده است: در حالت کوچک، خودرو پشت چراغ قرمز توقف می‌کند و وقتی چراغ سبز می‌شود شبه‌نام خود را تغییر می‌دهد، و در حالت بزرگ، خودرو در پارکینگ توقف می‌کند و هنگام خروج از پارکینگ شبه‌نام عوض می‌کند \cite{luPseudonymChangingSocial2012}. این نقاط به‌طور طبیعی به نواحی اختلاط تبدیل می‌شوند و شرط تمایزناپذیری برقرار است وقتی مکان برابر نقطهٔ اجتماعی و سرعت برابر صفر باشد \cite{luPseudonymChangingSocial2012}. معیار سنجش گمنامی در این طرح اندازهٔ مجموعهٔ گمنامی است؛ هرچه این اندازه بزرگ‌تر باشد، گمنامی بالاتر است \cite{luPseudonymChangingSocial2012}. برای نقطهٔ کوچک، ورود خودروها با فرایند پواسون و فاصلهٔ میانگین ورود $\lambda$ مدل می‌شود و با فرض زمان سکونت $t$، اندازهٔ مجموعهٔ گمنامی برابر $\lambda t$ است؛ برای نقطهٔ بزرگ، زمان از باز شدن مرکز خرید تا خروج خودرو نمایی با میانگین $\mu$ است و با در نظر گرفتن مدت ماندن هر خودرو با میانگین $\omega$، اندازهٔ مجموعهٔ گمنامی برابر $\lambda/(\omega+\mu)$ به دست می‌آید \cite{luPseudonymChangingSocial2012}.

مدل تهدید این طرح، مهاجمی سراسری و خارجی است که همهٔ پیام‌های ایمنی را با تجهیزات رادیویی و احتمالاً زیرساخت شنود جمع‌آوری می‌کند و ردیابی را به‌صورت مکانی-زمانی با زمان، مکان و سرعت انجام می‌دهد؛ این مهاجم صرفاً شنود منفعل است و به دنبال نفوذ به خودروها نیست، چون واحد درون‌خودرویی هنگام حرکت قابل خاموش‌کردن نیست و دوربین نیز به‌دلیل هزینهٔ بالا خارج از دامنه است \cite{luPseudonymChangingSocial2012}. انگیزهٔ بهره‌گیری از شبه‌نام آشکار است: با پنهان‌کردن هویت واقعی، حریم خصوصی هویتی به دست می‌آید، اما هویت توسط مرجع مورد اعتماد در صورت اختلاف قابل افشای شرطی است \cite{luPseudonymChangingSocial2012}. شبه‌نام ایستا کافی نیست، زیرا خودرو با اطلاعات مکانی-زمانی پیام‌ها قابل ردیابی است و پیوند بلندمدت — مانند شناسایی خانه یا محل کار راننده — هویت را لو می‌دهد \cite{luPseudonymChangingSocial2012}. تغییر سادهٔ شبه‌نام نیز کافی نیست؛ مهاجم می‌تواند از تکنیک‌های ردیابی چندهدفه برای پیوند دادن پیام‌های ارسال‌شده با شبه‌نام‌های مختلف استفاده کند و در حملهٔ نحوی، تنها خودرویی که در یک بازهٔ زمانی کوتاه شبه‌نام خود را عوض کرده شناسایی می‌شود و در حملهٔ معنایی، مسیر خودرو با مسیر همسایگان تفاوت دارد و مهاجم با ردیابی دو شبه‌نام را پیوند می‌زند؛ نتیجه این است که شبه‌نام‌ها تنها باید در موقعیت‌های مشاهده‌نشده تغییر کنند \cite{luPseudonymChangingSocial2012}.

دو رویکرد اصلی برای رسیدن به «موقعیت مشاهده‌نشده» وجود دارد. نخست، ناحیهٔ اختلاط: تغییر شبه‌نام در نواحی از پیش تعیین‌شده مانند تقاطع‌ها، نقاط اجتماعی و پمپ‌بنزین، که نیازمند زیرساختی برای اعلام مرز ناحیه است، پرهزینه و نایعمل به نظر می‌رسد و در برابر حملات زمان‌بندی و گذار آسیب‌پذیر است، زیرا مهاجم با پایش نقاط ورود و خروج و زمان ماندن، شبه‌نام قدیم و جدید را پیوند می‌زند \cite{luPseudonymChangingSocial2012}. دوم، دورهٔ سکوت: خودرو بدون نیاز به زیرساخت خودش تصمیم می‌گیرد مدتی پیام نفرستد و سپس با شبه‌نام جدید ادامه دهد؛ این تصمیم یا بر اساس زمان یا بر اساس بافت است و بیشتر پژوهشگران و استانداردها دورهٔ سکوت را چون نیازمند زیرساخت نبودن ترجیح داده‌اند \cite{luPseudonymChangingSocial2012}.

برای مقابله با خطر سرقت خودرو — با بیش از ۱۷۰ هزار خودروی دزدیده‌شده در سال در کانادا — طرح خودتفویضی شبه‌نام با کلید ایزوله \LTRfootnote{\lr{Key-insulated Pseudonym self-delegation (KPSD)}} را پیشنهاد می‌کند \cite{luPseudonymChangingSocial2012}. در این طرح، مرجع مورد اعتماد کلید ناشناس مجاز را به کاربر — نه خودرو — می‌دهد؛ کاربر آن را در محیط امنی مانند خانه نگه می‌دارد و پیش از سفر کلیدهای کوتاه‌عمر خودتفویض‌شده تولید و در واحد درون‌خودرویی نصب می‌کند، که قیاس آن به سوخت‌رسانی خودرو پیش از سفر شده است \cite{luPseudonymChangingSocial2012}. این طرح بر امضای کوتاه بونه-بوین و طرح \lr{ECPP} بنا شده و چهار مرحله دارد: راه‌اندازی سامانه، تولید کلید، تولید خودتفویض‌شدهٔ شبه‌نام و ردیابی شرطی \cite{luPseudonymChangingSocial2012}. از نظر کارایی، هزینهٔ وارسی $n$ پیام از یک منبع در این طرح $(3+n)T_{\mathrm{pair}} + (4+n)T_{\mathrm{exp1}} + 5T_{\mathrm{exp2}}$ است در حالی که طرح گروهی خالص \lr{GSB} برابر $3nT_{\mathrm{pair}} + 4nT_{\mathrm{exp1}} + 5nT_{\mathrm{exp2}}$ است و با فرض $T_{\mathrm{pair}} = 4.5$ میلی‌ثانیه، وارسی در محدودیت زمانی ۳۰۰ میلی‌ثانیه نمونه‌گیری شده است \cite{luPseudonymChangingSocial2012}.

آنچه این طرح را متمایز می‌کند، تحلیل نظریهٔ بازی آن است \cite{luPseudonymChangingSocial2012}. در مدل با $N$ خودرو، هر خودرو یا شبه‌نام خود را تغییر می‌دهد یا آن را نگه می‌دارد؛ هر خودرو ارزش‌گذاری شخصی $d_j$ برای حریم خصوصی و هزینهٔ نرمال‌شدهٔ $c_j$ برای تغییر دارد. با نگه‌داشتن، احتمال ردیابی یک است و پرداخت $-d_j$ است؛ با تغییر، احتمال ردیابی $1/S$ است — با کران پایین $S = n_{\mathrm{pm}}+1$ که $n_{\mathrm{pm}}$ تعداد همسایگانی است که شبه‌نام عوض کرده‌اند — و پرداخت $-d_j/(n_{\mathrm{pm}}+1) - c_j$ است \cite{luPseudonymChangingSocial2012}. شرط تغییر دادن برای یک خودرو $c_j < n_{\mathrm{pm}}{\cdot}d_j/(n_{\mathrm{pm}}+1)$ است و اگر هیچ همسایه‌ای شبه‌نام عوض نکرده باشد، این شرط برقرار نیست \cite{luPseudonymChangingSocial2012}. بهرهٔ حریم خصوصی مکانی $(n_{\mathrm{pm}}/(n_{\mathrm{pm}}+1)){\cdot}d_j$ است که در مقدار تغییر صعودی است و بیشینه آن وقتی که همهٔ خودروها شبه‌نام عوض کنند، برابر $((N-1)/N){\cdot}d_j$ است — وضعیتی که «برای همه سودمند» توصیف شده است \cite{luPseudonymChangingSocial2012}. چون این طرح شبه‌نام‌های کافی و ارزان فراهم می‌کند، هزینهٔ تغییر می‌تواند بسیار پایین باشد \cite{luPseudonymChangingSocial2012}. در پایان، جایگاه نوآورانهٔ این کار در آن است که پیشینیان عمدتاً صرفاً شبیه‌سازی کرده‌اند و این طرح مدل تحلیلی ارائه می‌دهد؛ آنتروپی تغییر شبه‌نام نیز با $H(PC) = -\sum P_{i \to PC} \log_2 P_{i \to PC}$ سنجیده می‌شود و در توزیع یکنواخت بیشینه برابر $\log_2 N$ است \cite{luPseudonymChangingSocial2012}.

\subsubsection{طرح حریم خصوصی مرتبط با ایمنی \lr{(SRPS)}}

در مقابل رویکرد پیشین که بر گمنامی در نقاط اجتماعی تمرکز دارد، طرح حریم خصوصی مرتبط با ایمنی \LTRfootnote{\lr{Safety-related Privacy Scheme (SRPS)}} مسئله‌ای دیگر را هدف می‌گیرد: پیام‌های بیکن در کاربردهای ایمنی، موقعیت، سرعت و جهت خودرو را به‌صورت متن ساده و با نرخ ۱ تا ۱۰ هرتز منتشر می‌کنند و شنودگر می‌تواند با پیوند دادن پیام‌های بیکن پیاپی، خودرو را ردیابی کند \cite{al-aniPrivacySafetyImprovement2023}. طرح‌های مبتنی بر دورهٔ سکوت حریم خصوصی را بهتر می‌کنند، اما در این دوره خودرو وضعیت خود را منتشر نمی‌کند و «تصادف احتمالی در این دوره قابل جلوگیری نیست» و تعادل بهینه میان حریم خصوصی و ایمنی همچنان چالشی باز است \cite{al-aniPrivacySafetyImprovement2023}.

طرح \lr{SRPS} با دو الگوریتم — \lr{SRPS-Active} و \lr{SRPS-Silent} — این تعادل را جست‌وجو می‌کند \cite{al-aniPrivacySafetyImprovement2023}. هر خودرو یک ردیاب چندهدفه \LTRfootnote{\lr{Multi-Target Tracker (MTT)}} دارد که با آبکرد کالمن، گیتینگ، انتساب دادهٔ \lr{NNPDA} و نگهداری ردیاب، وضعیت همسایگان را حتی در دورهٔ سکوت پیش‌بینی می‌کند؛ خودروی ساکت اگر برای گام بعد تصادفی پیش‌بینی کند از سکوت خارج می‌شود و خودروی فعال اگر وضعیتش با همسایه‌ای قابل اختلاط باشد، بدون ورود به سکوت شبه‌نام عوض می‌کند \cite{al-aniPrivacySafetyImprovement2023}. انگیزهٔ این طرح آشکار است: ارتباط کاملاً ناشناس پذیرفتنی نیست، زیرا کاربردهای ایمنی حیاتی‌اند و پاسخ‌گویی لازم است؛ بنابراین به‌جای هویت واقعی از شبه‌نام استفاده می‌شود تا میان امنیت و حریم خصوصی تعادل برقرار شود و شبه‌نام باید توسط طرف مورد اعتمادی صادر شود که بتواند در صورت اختلاف آن را حل کند \cite{al-aniPrivacySafetyImprovement2023}. آنچه تغییر سادهٔ شبه‌نام را ناکافی می‌کند، دو حمله است: در حملهٔ نحوی، خودرو تنها خودرویی است که در یک بازهٔ زمانی کوتاه شبه‌نامش را عوض کرده و در حملهٔ معنایی، مسیر خودرو با مسیر همسایگان تفاوت دارد و مهاجم با ردیابی دو شبه‌نام را پیوند می‌زند؛ نتیجه این است که شبه‌نام‌ها تنها باید در موقعیت‌های مشاهده‌نشده تغییر کنند \cite{al-aniPrivacySafetyImprovement2023}. هزینهٔ تغییر بیشتر شبه‌نام و سکوت طولانی‌تر بر ایمنی نیز دوگانه است: افزایش سربار امنیتی و از دست رفتن پیام‌ها — چون با شبه‌نام جدید یا همسایهٔ جدید گواهی باید پیوست و تأیید شود — و احتمال وقوع تصادف در دوره‌های سکوت که خودرو از اشتراک‌گذاری موقعیت خود دست می‌کشد \cite{al-aniPrivacySafetyImprovement2023}. از این رو، طراحی طرحی که سه مسئلهٔ کلیدی حریم خصوصی، امنیت و ایمنی را به‌طور مؤثر برآورده کند، همچنان چالشی علمی است \cite{al-aniPrivacySafetyImprovement2023}.

در ادبیات، طرح‌های پیشین شامل به‌روزرسانی دوره‌ای با دورهٔ ثابت — که برای مهاجم قابل پیش‌بینی است — یا تصادفی، تغییر همکارانه، ناحیهٔ اختلاط و دورهٔ سکوت تصادفی بوده‌اند که مشکل آن آخری این است که اگر تنها یک خودرو در جاده باشد، با وجود سکوت قابل شناسایی است \cite{al-aniPrivacySafetyImprovement2023}. طرح \lr{SLOW} تنها زیر آستانهٔ سرعت ساکت می‌شود که احتمال تصادف کمتری دارد، و طرح \lr{CAPS} خودروها را همکارانه وارد سکوت می‌کند و زمانی که بافت آن‌ها احتمالاً با خودروی ساکت دیگری اختلاط پیدا کند یا در موقعیت غیرمنتظره باشند، از سکوت خارج می‌شود و از ردیاب چندهدفه برای پیش‌بینی وضعیت خودروی ساکت استفاده می‌کند \cite{al-aniPrivacySafetyImprovement2023}. شکاف این پیشینیان آشکار است: تعداد کمی از طرح‌ها تأثیر بر کاربردهای ایمنی را در نظر گرفته‌اند و هیچ‌کدام تصادفات احتمالی در دوره‌های سکوت را — که انگیزهٔ این کار بود — برطرف نکرده‌اند \cite{al-aniPrivacySafetyImprovement2023}.

مدل سامانه شامل واحد درون‌خودرویی در هر خودرو است که پیام بیکن شامل موقعیت، سرعت و جهت را با نرخ ۱ تا ۱۰ هرتز در برد ۳۰۰ متر از طریق ارتباطات اختصاصی کوتاه‌برد ارسال می‌کند؛ واحدهای کنارجاده‌ای در کنار جاده نصب‌اند و ارتباط آن‌ها با یکدیگر و با مرجع سیمی و با خودروها بی‌سیمی است \cite{al-aniPrivacySafetyImprovement2023}. در زیرساخت کلید عمومی، کلید عمومی تأییدشده بدون اطلاعات هویتی به‌عنوان شبه‌نام عمل می‌کند و در دستگاه مقاوم در برابر دست‌کاری نگهداری می‌شود و نقش‌های مراجع تفکیک شده‌اند \cite{al-aniPrivacySafetyImprovement2023}. مدل مهاجم در این طرح، مهاجم سراسری منفعل است که با شنود و پایش همهٔ پیام‌های پخش‌شده به حریم خصوصی خودروها هجوم می‌برد؛ نمونه‌ای از آن ارائه‌دهندهٔ سرویس نامطمئن است \cite{al-aniPrivacySafetyImprovement2023}.

مدیریت شبه‌نام در این طرح با دو مرجع مورد اعتماد — مرجع صدور بلندمدت و مرجع صدور کوتاه‌مدت — که هرکدام با یک زوج کلید مجهز‌اند، انجام می‌شود \cite{al-aniPrivacySafetyImprovement2023}. خودرو با ارائهٔ مدارک، شبه‌نام بلندمدت را از مرجع بلندمدت دریافت می‌کند؛ این شبه‌نام هویت ایستای خودرو است، تنها با تغییر مالک عوض می‌شود و با کلید خصوصی مرجع بلندمدت امضا می‌شود \cite{al-aniPrivacySafetyImprovement2023}. خودرو درخواست شبه‌نام کوتاه‌مدت را با کلید خصوصی متناظر شبه‌نام بلندمدت — مستقیماً یا با کمک واحد کنارجاده‌ای — امضا می‌کند؛ مرجع کوتاه‌مدت اعتبار شبه‌نام بلندمدت را با کلید عمومی مرجع بلندمدت و اصالت درخواست را با کلید عمومی داخل گواهی بررسی می‌کند و شبه‌نام کوتاه‌مدت صادر می‌کند \cite{al-aniPrivacySafetyImprovement2023}. شبه‌نام کوتاه‌مدت برای احراز اصالت پیام‌های ایمنی به‌کار می‌رود: ابتدا مهر زمانی افزوده و سپس با کلید خصوصی شبه‌نام کوتاه‌مدت فعلی امضا می‌شود و مهر زمانی از پخش مجدد جلوگیری می‌کند — مانند پخش مجدد پیام خودروی امدادی \cite{al-aniPrivacySafetyImprovement2023}. هر شبه‌نام کوتاه‌مدت هم حداقل عمر — برای پیوندناپذیری کوتاه‌مدت — و حداکثر عمر — برای جلوگیری از پیوند بلندمدت — دارد و شبه‌نام‌های بلندمدت و کوتاه‌مدت در دستگاه مقاوم در برابر دست‌کاری نگهداری می‌شوند و مراجع صادرکننده پایگاه‌داده‌ای از پیوند هویت واقعی به شبه‌نام بلندمدت و آن به شبه‌نام‌های کوتاه‌مدت برای پاسخ‌گویی نگه می‌دارند \cite{al-aniPrivacySafetyImprovement2023}.

ردیاب خودرو با چهار مرحله، هم ابزار مهاجم برای پیوند شبه‌نام‌ها و هم دفاع درون خودرو را تجسم می‌کند \cite{al-aniPrivacySafetyImprovement2023}: تخمین وضعیت با آبکرد کالمن که اندازه‌گیری‌های نادقیق حسگر و تخمین مدل سینماتیکی را ترکیب می‌کند و بهترین تخمین موقعیت، سرعت و جهت را می‌دهد؛ گیتینگ برای حذف انتساب‌های نامحتمل پیش از انتساب؛ انتساب داده برای پیوند پیام‌های یک خودرو با شبه‌نام‌های مختلف از طریق ماتریس احتمال انتساب با روش \lr{NNPDA} که محاسبهٔ بلادرنگ را حتی با تعداد زیاد خودرو ممکن می‌کند؛ و نگهداری ردیاب برای حذف خودروهای خارج از برد و ادامهٔ ردیابی همسایگان حتی اگر ساکت باشند \cite{al-aniPrivacySafetyImprovement2023}. همان ابزاری که مهاجم برای پیوند شبه‌نام‌ها به کار می‌برد، در این طرح به‌صورت دفاعی در خود خودرو اجرا می‌شود تا سطح سردرگمی مهاجم را بسنجد و تصادف را پیش‌بینی کند \cite{al-aniPrivacySafetyImprovement2023}. تشخیص خودروی ساکت نیز زمانی صورت می‌گیرد که دو پیام پیاپی از همسایه دریافت نشود، با این توجه که یک پیام گم‌شده ممکن است ناشی از سربار باشد \cite{al-aniPrivacySafetyImprovement2023}.

اهداف طرح چهارگانه‌اند \cite{al-aniPrivacySafetyImprovement2023}. نخست، کاهش تصادف در دورهٔ سکوت: هر خودرو — ساکت یا فعال — در هر گام موقعیت بعدی خود و همسایگان را با آبکرد کالمن پیش‌بینی می‌کند و اگر خودروی ساکت برای گام بعد تصادفی پیش‌بینی کند، از دورهٔ سکوت خارج شده و انتشار وضعیت خود را آغاز می‌کند \cite{al-aniPrivacySafetyImprovement2023}. دوم، کاهش هدر رفتن شبه‌نام در موقعیت‌های مشاهده‌شده. سوم، کاهش دوره‌های سکوت: خودرو می‌تواند بدون ورود به دورهٔ سکوت شبه‌نام خود را با موفقیت تغییر دهد، به‌شرطی که وضعیتش احتمالاً با دیگران اختلاط پیدا کند \cite{al-aniPrivacySafetyImprovement2023}. چهارم، افزایش احتمال بافت اختلاط: خودروها همکارانه ساکت شده و بلافاصله به جست‌وجوی بافت اختلاط می‌پردازند؛ در مقابل، طرح \lr{CAPS} پس از ۳ ثانیه به جست‌وجو شروع می‌کند، در حالی که خودرو در ۳ ثانیه تا ۶۰ متر جابه‌جا می‌شود و فرصت اختلاط کمتر است \cite{al-aniPrivacySafetyImprovement2023}.

در الگوریتم \lr{SRPS-Active}، خودرو تا زمانی که طول دورهٔ فعالیت از حداقل طول — که ۶۰ ثانیه پیشنهاد شده — کمتر است، وضعیت خود را با شبه‌نام فعلی منتشر می‌کند \cite{al-aniPrivacySafetyImprovement2023}. در بازهٔ میان حداقل و حداکثر طول، سه شرط به ترتیب بررسی می‌شوند: وضعیت غیرمنتظره — اگر فاصلهٔ میان وضعیت پیش‌بینی‌شدهٔ گام قبل و وضعیت واقعی برای سردرگمی مهاجم کافی باشد، شبه‌نام تغییر و وضعیت منتشر می‌شود، مانند خودرویی که قصد مستقیم‌روی داشت اما به دلیل هشدار تصادف یا ترافیک می‌پیچد یا می‌ایستد؛ بافت اختلاط با همسایه — با به‌روزرسانی و پیش‌بینی کالمن برای وضعیت پیش‌بینی‌شده و پیام‌های دریافتی همهٔ همسایگان، ساکت و فعال، اگر وضعیت پیش‌بینی‌شده با یکی از پیام‌های همسایه برابر باشد، وضعیت فعلی منتشر و شبه‌نام برای وضعیت بعد تغییر می‌کند؛ و تغییر وضعیت همسایه — اگر همسایه‌ای ساکت شده باشد، ورود همکارانه به دورهٔ سکوت و فراخوانی الگوریتم \lr{SRPS-Silent} \cite{al-aniPrivacySafetyImprovement2023}. در غیر این صورت انتشار با همان شبه‌نام ادامه می‌یابد و اگر طول دوره از حداکثر طول بگذرد، سکوت اجباری برای جلوگیری از پیوند بلندمدت اعمال می‌شود \cite{al-aniPrivacySafetyImprovement2023}. در الگوریتم \lr{SRPS-Silent}، خودروی ساکت بلافاصله به جست‌وجوی فرصت ازسرگیری می‌پردازد: اگر وضعیت پیش‌بینی‌شده با وضعیت واقعی تفاوت داشته باشد، فعال می‌شود و شبه‌نام عوض می‌کند؛ اگر پیش‌بینی وضعیت و پیام‌های همسایگان حاکی از اختلاط بافت یا تصادف پیش‌بینی‌شده باشد، شمارندهٔ تصادفات پیش‌بینی‌شده افزایش، شبه‌نام تغییر و خودرو فعال می‌شود — نکتهٔ کلیدی آنکه همان شرطی که فرصت حریم خصوصی است، هشدار ایمنی نیز هست، زیرا اختلاط بافت یعنی دو خودرو احتمالاً در یک موقعیت یا نزدیک یکدیگر قرار می‌گیرند که اگر خودرو انتشار وضعیت خود را ادامه ندهد، می‌تواند به تصادف منجر شود — و اگر دورهٔ سکوت از حداکثر آن بگذرد، خروج اجباری از سکوت اعمال می‌شود \cite{al-aniPrivacySafetyImprovement2023}.

ارزیابی این طرح در محیط \lr{OMNeT++} نسخهٔ ۵.۰ و \lr{SUMO} نسخهٔ ۰.۲۵.۰ و \lr{Veins} نسخهٔ ۴.۴ با اتصال \lr{PREXT} به \lr{TraCI} انجام شده است، بر روی نقشهٔ ۳.۸ در ۲.۸ کیلومتری از لیورپول بریتانیا با نرخ ورود یک خودرو در هر ۱، ۰.۵ و ۰.۳ ثانیه، هر آزمون ۳۶۰ ثانیه‌ای و میانگین سه پایگاه سفر تصادفی و نرخ بیکن ۱۰ هرتز \cite{al-aniPrivacySafetyImprovement2023}. طرح \lr{SRPS} با پنج طرح \lr{PPC}، \lr{RSP}، \lr{CSP}، \lr{SLOW} و \lr{CAPS} مقایسه شده و پارامترهای آن عبارت‌اند از حداقل طول فعالیت ۶۰ ثانیه، حداکثر طول فعالیت ۱۲۰ ثانیه، حداقل سکوت صفر ثانیه، حداکثر سکوت ۱۳ ثانیه و شعاع همسایگی ۵۰ متر \cite{al-aniPrivacySafetyImprovement2023}. معیارهای سنجش شامل سربار امنیتی (تعداد تغییر شبه‌نام در ثانیه)، ردیابی‌پذیری (درصد خودروهایی که برای حداقل ۹۰ درصد عمرشان به‌طور پیوسته قابل ردیابی‌اند و خودروهایی که هرگز شبه‌نام عوض نکرده‌اند حذف می‌شوند)، ایمنی (پیام بیکن در ثانیه و تصادفات پیش‌بینی‌شده/اجتناب‌شده) و کارایی (درصد سردرگمی و تعداد خودروهایی که با وجود تغییر شبه‌نام قابل ردیابی ماندند) است \cite{al-aniPrivacySafetyImprovement2023}.

یافته‌ها — که مقادیر آن‌ها از برچسب‌های نمودارها خوانده شده و در متن مقاله نیز همین‌گونه گزارش شده — نشان می‌دهد ردیابی‌پذیری در \lr{SRPS} به‌ترتیب ۴۰، ۳۱ و ۲۰ درصد در سه نرخ ورود است در حالی که در \lr{CAPS} ۷۱، ۶۳ و ۵۲ درصد و در \lr{PPC} تا ۹۴ درصد است و \lr{CSP} حدود ۵ تا ۷ درصد نوسان دارد؛ متن مقاله از کاهش تقریبی ۳۰ درصدی ردیابی‌پذیری نسبت به \lr{CAPS} سخن می‌گوید \cite{al-aniPrivacySafetyImprovement2023}. تعداد پیام بیکن در ثانیه در \lr{SRPS} به‌ترتیب ۹.۸۱، ۹.۸۷ و ۹.۹۲ است و همواره بالای ۹ قرار دارد، در حالی که \lr{SLOW} همواره زیر ۶.۵۰ است و در میانگین ۳.۵ پیام در ثانیه از دست می‌دهد و \lr{RSP} زیر ۷.۶۵ است؛ مقاله از \lr{SRPS} به‌عنوان «تنها طرح حریم خصوصی که ارسال پیام را با افزایش تعداد خودروها کاهش نمی‌دهد» یاد می‌کند \cite{al-aniPrivacySafetyImprovement2023}. درصد سردرگمی در \lr{SRPS} ۶۹، ۷۷ و ۸۴ درصد است در برابر ۲۶، ۳۸ و ۴۱ درصد در \lr{CAPS} و مقاله افزایش «بیش از ۳۹ درصدی» سردرگمی را گزارش می‌کند، در حالی که \lr{CSP} با سردرگمی ۱۰۰ درصدی و کمترین هدر رفتن شبه‌نام (کمتر از ۱۹ خودرو)، ایمنی را در سکوت همگانی قربانی می‌کند و بیش از ۸۵ درصد خودروها هر دقیقه شبه‌نام عوض می‌کنند و «شبکه اقتضایی خودرویی وقتی همهٔ خودروها هم‌زمان وارد دورهٔ سکوت شوند، از کار می‌افتد» \cite{al-aniPrivacySafetyImprovement2023}. تصادفات پیش‌بینی‌شده/اجتناب‌شده در \lr{SRPS} به‌ترتیب ۴۳، ۶۸ و ۹۳ با افزایش تراکم گزارش شده و تعداد تغییر شبه‌نام در ثانیه ۰.۶۲، ۰.۶۶ و ۰.۷۳ است که با تراکم به‌صورت سازگار رشد می‌کند، در حالی که \lr{PPC} با ۰.۷۴ بالاترین است چون سکوتی ندارد \cite{al-aniPrivacySafetyImprovement2023}. در جدول رتبه‌بندی کلی، \lr{SRPS} در حریم خصوصی رتبهٔ ۳، در ایمنی رتبهٔ ۲ و در سربار رتبهٔ ۲ را دارد و نویسندگان از دستیابی به «بهترین تعادل» میان سه مسئلهٔ حریم خصوصی، ایمنی و کارایی سخن می‌گویند — ادعایی که به خود نویسندگان نسبت داده می‌شود \cite{al-aniPrivacySafetyImprovement2023}.

محدودیت‌های این طرح نیز باید در نظر گرفته شود. شعاع ثابت ۵۰ متر فقط برای جاده‌های شهری با حداکثر سرعت ۶۴ کیلومتر بر ساعت برگزیده شده و تنظیم پویا بر اساس تراکم، کار آینده است \cite{al-aniPrivacySafetyImprovement2023}. اگر همسایه شبه‌نام خود را عوض نکند، خودرو همچنان با اطلاعات مکانی-زمانی قابل پیوند است و این موضوع از دامنهٔ طرح خارج شده است \cite{al-aniPrivacySafetyImprovement2023}. ارزیابی صرفاً شبیه‌سازی روی یک نقشهٔ شهری است که عمداً برای افزایش احتمال بافت اختلاط برگزیده شده و تعداد تصادفات پیش‌بینی‌شده، شمار «تصادفات مورد انتظار» است نه تصادفات واقعاً شبیه‌سازی‌شده، و مدل مهاجم نیز صرفاً سراسری منفعل است \cite{al-aniPrivacySafetyImprovement2023}.

\subsubsection{سامانهٔ شناسه‌های غیرمتمرکز مبتنی بر شهرت \lr{(DIVA)}}

سامانهٔ \lr{DIVA} ترکیبی از شناسه‌های غیرمتمرکز و \lr{IOTA Tangle} برای مدیریت شهرت در شبکه اقتضایی خودرویی است و هدف آن انتقال امن پیام‌ها با حفظ حریم خصوصی فرستنده است \cite{feraudoDIVADIDbasedReputation2024}. در معماری این سامانه، مرجع مورد اعتماد شناسه‌های غیرمتمرکز را ثبت و اعتبارنامه صادر می‌کند و سه گرهٔ لبه — گرهٔ دفتر کل توزیع‌شده، مشتری دفتر کل توزیع‌شده و سامانهٔ شهرت — در هستهٔ نسل پنجم با عملکرد گرهٔ کاربرابری و عملکرد کاربرابری بین‌عملکردی و یک واحد کنارجاده‌ای دارای حسگر و یک واحد کنارجاده‌ای معمولی حضور دارند \cite{feraudoDIVADIDbasedReputation2024}.

\begin{figure}[H]
\centering
\includegraphics[width=0.9\textwidth]{DIVA.jpg}
\caption{معماری سیستم \lr{DIVA}: سیستم شهرت مبتنی بر شناسه‌های غیرمتمرکز برای انتقال امن در شبکه اقتضایی خودرویی با استفاده از \lr{IOTA Tangle}}
\label{fig:diva_arch}
\end{figure}

شکل \ref{fig:diva_arch} مدل سیستم \lr{DIVA} را نشان می‌دهد و نقش مرجع مورد اعتماد، گره‌های لبه و واحدهای کنارجاده‌ای را در حساب شهرت فرستندهٔ پیام ترسیم می‌کند.

در این سامانه، شناسهٔ غیرمتمرکز برای شناسایی خودرو به‌کار می‌رود و به‌عنوان شبه‌نام عمل می‌کند و قابل نسبت دادن به هویت واقعی کاربر نیست \cite{feraudoDIVADIDbasedReputation2024}. امتیازات شهرت به‌صورت تراکنش‌های بدون ارزش در یک دفتر کل با دسترسی محدود ذخیره می‌شوند \cite{feraudoDIVADIDbasedReputation2024}. امتیاز فرستندهٔ پیام اعلان محیطی غیرمتمرکز \LTRfootnote{\lr{Decentralized Environmental Notification Message (DENM)}} بر روی گرهٔ لبه محاسبه می‌شود، به‌این‌ترتیب که پیام فرستنده با دادهٔ واحد کنارجاده‌ای، پیام‌های آگاهی همیارانه \LTRfootnote{\lr{Cooperative Awareness Message (CAM)}} همان فرستنده و پیام‌های مشابه اعلان محیطی مقایسه می‌شود \cite{feraudoDIVADIDbasedReputation2024}. یکپارچگی و انکارناپذیری نیز با امضای پیام‌ها با کلید خصوصی و قرار دادن کلید عمومی در سند شناسه بر تانگل تضمین می‌شود و رمزنگاری مبنای طرح بر رمزنگاری خم بیضوی استوار است \cite{feraudoDIVADIDbasedReputation2024}.

در برابر حملهٔ سیبل، هر خودرو یک شناسهٔ غیرمتمرکز دارد و با سقوط شهرت به زیر آستانه، پنجرهٔ حمله محدود می‌شود، و در برابر جعل، رمزنگاری خم بیضوی و کلید خصوصی دفاع است \cite{feraudoDIVADIDbasedReputation2024}. ارزیابی سامانه بر مجموعهٔ دادهٔ \lr{V2V} اتحادیهٔ اروپایی برای ارتباطات خودرویی انجام شده و کد و داده در دسترس عمومی قرار دارند \cite{feraudoDIVADIDbasedReputation2024}. در روش تزریق نویز، برای هر ستون یک توزیع نرمال با میانگین صفر و انحراف معیار متناسب با همان ستون ساخته می‌شود و هر مقدار مستقلاً تغییر می‌کند تا داده واقع‌گرایانه بماند؛ آستانهٔ حدت رویداد با مد، میانه یا میانگین روی دادهٔ سالم تعیین شده و آستانهٔ پیام آگاهی همیارانه ۶۰۰ ثانیه است و ضریب میان از صفر تا یک با گام ۰.۱ تغییر کرده است \cite{feraudoDIVADIDbasedReputation2024}.

نتایج این طرح باید با دقت خوانده شوند. نویسندگان در مقدمه از شناسایی دقیق مشارکت‌کنندگان مخرب «با دقت تقریبی ۹۰ درصد در سناریوهای مختلف» سخن می‌گوند، اما در بخش ارزیابی اعداد شرطی «حدود ۹۴ درصد با ۳۰ درصد و ۸۹ درصد با ۴۰ درصد منبع مخرب» با آستانهٔ میانگین و ضریب بیش از ۰.۳ گزارش می‌شوند و در نتیجه‌گیری حتی «حدود ۹۹ درصد با آستانه‌های مشخص» آمده که این ناسازگاری درونی منبع است \cite{feraudoDIVADIDbasedReputation2024}. این دقت‌ها بر پیام‌های اعلان محیطی در یک سناریوی شبیه‌سازی‌شده سنجیده شده‌اند و تأخیر «چند میکروثانیه» نیز مربوط به اندازه‌گیری شبیه‌سازی یک‌به‌یک با پنج گرهٔ لبه و دفتر کل با دسترسی محدود است \cite{feraudoDIVADIDbasedReputation2024}. یک شکاف پژوهشی نیز در این طرح وجود دارد: چون شناسه‌های غیرمتمرکز به هویت واقعی کاربر نمی‌توانند نسبت داده شوند، ردیابی شرطیِ هویت واقعی وجود ندارد و ردیابی تنها از طریق فیلتر کردن شهرت امکان‌پذیر است \cite{feraudoDIVADIDbasedReputation2024}.

در کارهای مرتبط، طرح‌هایی مانند \lr{BDRA} با بلاکچین دولایه و شناسهٔ غیرمتمرکز — که نحوهٔ بررسی درستی پیام‌ها در آن توضیح داده نشده — و \lr{BRS4VANET} با بلاکچین کنسرسیومی و قرارداد هوشمند — که سازوکار ابطال آن مشخص نیست — و \lr{BARS} با شبه‌نام مبتنی بر کلید عمومی و شهرت مستقیم و غیرمستقیم — که بررسی درستی پیام در آن نیامده — و طرح‌های مبتنی بر گراف جهت‌دار بی‌دور پارتیشن‌بندی‌شده و سازگاری محلی و کنترل نرخ مبتنی بر شهرت پیشنهاد شده‌اند؛ مزیت‌های گراف جهت‌دار بی‌دور نیز «مقیاس‌پذیری، تراکنش‌های سریع و کارمزد پایین» گزارش شده است \cite{feraudoDIVADIDbasedReputation2024}.

\subsubsection{پروتکل‌های امنیتی مبتنی بر زنجیره بلوکی}

پروتکل‌های مبتنی بر زنجیره بلوکی، وابستگی طرح‌های کلاسیک مبتنی بر امضای گروهی و امضای مبتنی بر هویت — مانند \lr{GSIS} که در فصل ۶ بررسی شد و برای ردیابی شرطی به یک مرجع متمرکز متکی است — را به سمت اعتماد توزیع‌شده می‌برند \cite{linGSISSecurePrivacyPreserving2007}. دو طرح نمایان در این دسته هستند.

نخست، سامانهٔ احراز هویت حافظ حریم خصوصی با کمک بلاکچین \LTRfootnote{\lr{Blockchain-Assisted Privacy-Preserving Authentication System (BPAS)}} چارچوبی است که احراز هویت را به‌صورت خودکار در شبکه‌های اقتضایی خودرویی انجام می‌دهد و هم‌زمان حریم خصوصی خودرو را حفظ می‌کند \cite{fengBPASBlockchainAssistedPrivacyPreserving2020}. بر اساس چکیدهٔ این مقاله، این سامانه نیازی به مرکز ثبت آنلاین ندارد — مگر در دو مرحلهٔ راه‌اندازی سامانه و ثبت‌نام خودرو — و امکان ردیابی شرطی و ابطال پویا خودروهای رفتار نادرست را فراهم می‌کند و نویسندگان آن را «بسیار کارا و مقیاس‌پذیر» می‌نامند \cite{fengBPASBlockchainAssistedPrivacyPreserving2020}. تحلیل امنیتی عمیق و ارزیابی عملکرد جامع این طرح بر بستر \lr{Hyperledger Fabric} — که یک زنجیره بلوکی با دسترسی محدود است — انجام شده و نتیجهٔ اعلامی نویسندگان، ارائهٔ «راه‌حلی کارا برای توسعهٔ سامانهٔ احراز هویت غیرمتمرکز در شبکه‌های اقتضایی خودرویی» است \cite{fengBPASBlockchainAssistedPrivacyPreserving2020}. باید توجه داشت که در مرور حاضر تنها چکیده و پاراگراف نخست مقدمهٔ این مقاله در دسترس بوده و جزئیات سازوکار قرارداد هوشمند، نهاد ردیاب و سازوکار ابطال در آن قابل بررسی نبوده است \cite{fengBPASBlockchainAssistedPrivacyPreserving2020}.

دوم، طرح احراز هویت ناشناس و حفظ یکپارچگی مبتنی بر بلاکچین \LTRfootnote{\lr{Blockchain-based Anonymous Authentication and Integrity Preservation (BAIV)}} مسئلهٔ احراز هویت مجدد خودرو هنگام عبور از ناحیهٔ یک واحد کنارجاده‌ای به ناحیهٔ واحد بعدی — که هزینهٔ محاسباتی را بالا می‌برد — را هدف می‌گیرد و با یکپارچ‌سازی بلاکچین با شبکه اقتضایی خودرویی، احراز هویت را بدون مشارکت مرجع مورد اعتماد ممکن می‌سازد و حریم خصوصی شرطی را برای ابطال خودروهای مخرب در صورت اختلاف معرفی می‌کند \cite{mariaBAIVEfficientBlockchainBased2022}. در مدل این طرح، سه موجودیت مرجع مورد اعتماد، واحد کنارجاده‌ای و واحد درون‌خودرویی به‌علاوهٔ یک شبکهٔ بلاکچین مستقل در نظر گرفته شده‌اند و همهٔ واحدهای کنارجاده‌ای و مرجع مورد اعتماد به‌طور مستقل به بلاکچین وصل‌اند؛ اطلاعات خودرو پس از نخستین احراز هویت در بلاکچین ذخیره می‌شود و واحدهای کنارجاده‌ای بعدی نیازی به احراز هویت مجدد ندارند \cite{mariaBAIVEfficientBlockchainBased2022}. هر بلوک از طریق هش بلوک قبلی به هم پیوند می‌خورد، اطلاعات به‌صورت هش \lr{SHA256} نمایش داده می‌شوند و چون هیچ طرف سومی بلاکچین را اداره نمی‌کند، سامانه کاملاً غیرمتمرکز است \cite{mariaBAIVEfficientBlockchainBased2022}. طرح در هشت مرحله — راه‌اندازی سامانه، ثبت‌نام خودرو، ثبت‌نام واحد کنارجاده‌ای، احراز هویت ناشناس کاربر خودرو، احراز هویت ناشناس واحد کنارجاده‌ای، تحویل هنگام دست‌به‌دست‌شدن با «رسید احراز هویت»، انتقال امن پیام و حفظ یکپارچگی، و ابطال — پیاده‌سازی شده و پایهٔ رمزنگاری آن رمزنگاری خم بیضوی با نگاشت دوخطی، تابع درهم‌ساز، عمل \lr{XOR} و سختی مسئلهٔ لگاریتم گسسته است \cite{mariaBAIVEfficientBlockchainBased2022}.

عملکرد این طرح در محیط سیگوین با هستهٔ \lr{i7} ۳.۴ گیگاهرتز و ۸ گیگابایت رم و کتابخانهٔ \lr{PBC} سنجیده شده و زمان عملیات پایه عبارت‌اند از جمع نقطه‌ای ۰.۰۱۱ میلی‌ثانیه، ضرب نقطه‌ای ۲.۴ میلی‌ثانیه، درهم‌سازی ۰.۰۱ میلی‌ثانیه، نگاشت دوخطی ۲.۹ میلی‌ثانیه و \lr{XOR} ۰.۰۱ میلی‌ثانیه \cite{mariaBAIVEfficientBlockchainBased2022}. هزینهٔ احراز هویت یک کاربر و یک واحد کنارجاده‌ای ۱۲.۵۴ میلی‌ثانیه است، یعنی تقریباً ۸۰ کاربر در ثانیه و «نزدیک به ۴۸۰۰» کاربر در دقیقه، و برای ۱۰۰ کاربر ۱۲۹۱ میلی‌ثانیه در مقابل «بیش از ۱۳۷۵ میلی‌ثانیه» برای طرح‌های دیگر گزارش شده که این مقادیر میانگین ۱۰۰ شبیه‌سازی تصادفی هستند \cite{mariaBAIVEfficientBlockchainBased2022}. تحلیل امنیتی این طرح صرفاً استدلالی است و در برابر جعل هویت، دستکاری پیام، پیام نادرست، یکپارچگی، پیوندناپذیری، پخش مجدد، ردیابی شرطی، حریم خصوصی شرطی و انکارناپذیری استدلال توصیفی می‌آورد و هیچ اثبات رسمی — نه کاهش در مدل مهاجم، نه تأیید با ابزارهایی مانند \lr{ProVerif} یا \lr{AVISPA} — ارائه نمی‌دهد \cite{mariaBAIVEfficientBlockchainBased2022}. محدودیت‌های آن نیز شامل نامشخص‌بودن نوع بلاکچین، سازوکار اجماع و اندازهٔ بلوک، عدم ارزیابی عملی بلاکچین در پیاده‌سازی (تنها عملیات رمزنگاری زمان‌سنجی شده)، عدم تحلیل هزینهٔ ارتباطی و نیاز همچنان به مرجع مورد اعتماد برای ثبت‌نام و ابطال است \cite{mariaBAIVEfficientBlockchainBased2022}.
